\documentclass[11pt]{article}
\usepackage[utf8]{inputenc}
\usepackage{amsmath}
\usepackage{listings}
\usepackage{xcolor}
\usepackage{booktabs}
\usepackage{geometry}

\geometry{a4paper, margin=1in}

\definecolor{codegreen}{rgb}{0,0.6,0}
\definecolor{codegray}{rgb}{0.5,0.5,0.5}
\definecolor{codepurple}{rgb}{0.58,0,0.82}
\definecolor{backcolour}{rgb}{0.95,0.95,0.92}

\lstset{
    backgroundcolor=\color{backcolour},   
    commentstyle=\color{codegreen},
    keywordstyle=\color{magenta},
    numberstyle=\tiny\color{codegray},
    stringstyle=\color{codepurple},
    basicstyle=\ttfamily\small,
    breaklines=true,                 
    captionpos=b,                    
    keepspaces=true,                 
    numbers=left,                    
    numbersep=5pt,                  
    tabsize=2
}

\begin{document}

\title{Documentation: Multiplexer Logic Realization (Q.39)}
\author{Digital Systems Analysis}
\date{}
\maketitle

\section{Step 1: The Question}
Determine the Boolean function $F(A, B, C, D)$ realized by the $4 \times 1$ multiplexer circuit shown in the diagram.

\textbf{Options:}
\begin{itemize}
    \item (A) $F = \sum m(0, 1, 3, 5, 9, 10, 14)$
    \item (B) $F = \sum m(2, 3, 5, 7, 8, 12, 13)$
    \item (C) $F = \sum m(1, 2, 4, 5, 11, 14, 15)$
    \item (D) $F = \sum m(2, 3, 5, 7, 8, 9, 12)$
\end{itemize}

---

\section{Step 2: Logic Analysis}
A $4 \times 1$ MUX selects one of four data inputs based on select lines $S_1$ and $S_0$. The general equation is:
\[ F = \bar{S_1}\bar{S_0}I_0 + \bar{S_1}S_0I_1 + S_1\bar{S_0}I_2 + S_1S_0I_3 \]

\textbf{Circuit Connections:}
\begin{itemize}
    \item \textbf{Select Lines:} $S_1 = A$, $S_0 = B$.
    \item \textbf{Input $I_0$:} Connected to $C$.
    \item \textbf{Input $I_1$:} Connected to $D$.
    \item \textbf{Input $I_2$:} Connected to $\bar{C}$ (via NOT gate).
    \item \textbf{Input $I_3$:} Connected to $\bar{C}\bar{D}$ (via bubbled AND gate).
\end{itemize}

\textbf{Minterm Expansion:}
\begin{enumerate}
    \item $\bar{A}\bar{B}C = \bar{A}\bar{B}C(D + \bar{D}) = \bar{A}\bar{B}CD + \bar{A}\bar{B}C\bar{D} \rightarrow m_3, m_2$.
    \item $\bar{A}BD = \bar{A}B(C + \bar{C})D = \bar{A}BCD + \bar{A}B\bar{C}D \rightarrow m_7, m_5$.
    \item $A\bar{B}\bar{C} = A\bar{B}\bar{C}(D + \bar{D}) = A\bar{B}\bar{C}D + A\bar{B}\bar{C}\bar{D} \rightarrow m_9, m_8$.
    \item $AB\bar{C}\bar{D} \rightarrow m_{12}$.
\end{enumerate}
Combined, $F = \sum m(2, 3, 5, 7, 8, 9, 12)$.

---

\section{Step 3: Python Code Implementation}
The following script evaluates all 16 input combinations to identify the active minterms.

\begin{lstlisting}[language=Python]
import numpy as np

# Generate all 16 combinations for A, B, C, D (0-15)
inputs = np.array([[(i>>3)&1, (i>>2)&1, (i>>1)&1, i&1] for i in range(16)])
A, B, C, D = inputs[:,0], inputs[:,1], inputs[:,2], inputs[:,3]

def solve_mux():
    # Define MUX inputs from circuit
    I0, I1 = C, D
    I2, I3 = (1 - C), (1 - C) & (1 - D)
    
    # MUX logic realization
    F = ((1-A) & (1-B) & I0) | ((1-A) & B & I1) | \
        (A & (1-B) & I2) | (A & B & I3)
    
    # Extract indices where output is High
    minterms = np.where(F == 1)[0]
    print(f"Minterms found: {list(minterms)}")

if __name__ == "__main__":
    solve_mux()
\end{lstlisting}

---

\section{Step 4: Explanation of the Code}
\begin{itemize}
    \item \textbf{Binary Generation:} We use bit-shifting to create a $16 \times 4$ matrix representing the truth table inputs for $A, B, C, D$.
    \item \textbf{Gate Modeling:} The NOT gate is modeled as \texttt{1 - C}, and the bubbled AND gate is modeled using the bitwise \texttt{\&} of two inverted signals.
    \item \textbf{Logical Sum:} The final expression aggregates the selection logic for all four MUX channels.
\end{itemize}

---

\section{Step 5: Output and Conclusion}
The script identifies the minterms as follows:
\begin{lstlisting}[backgroundcolor=\color{black}, basicstyle=\ttfamily\color{white}]
--- Truth Table Minterms ---
Minterms found: [2, 3, 5, 7, 8, 9, 12]
\end{lstlisting}

\textbf{Final Result:} The realized function matches \textbf{Option (D)}.

\end{document}
